%! AP Statistics — Unit 9, Lesson 9.5 — Homework (Answer Key)
\documentclass[10pt]{article}
\usepackage{apstats-article}
\usepackage{apstats-key}

\begin{document}

\pageheader{Unit 9, Lesson 9.5}{Homework --- Answer Key}
\namedateperiod

\vspace{0.1in}
\textbf{Directions:} For each problem, carry out a significance test for slope as directed.
Show all work clearly.

\begin{enumerate}

  \item \textbf{Age and Running Speed.} A sports scientist studies the relationship between
        age (years, $x$) and maximum running speed (mph, $y$) in adult recreational runners
        (random sample, $n = 28$).

        \begin{center}
        \renewcommand{\arraystretch}{1.4}
        \begin{tabular}{lrrrr}
            \textbf{Predictor} & \textbf{Coef} & \textbf{SE Coef} & \textbf{T} & \textbf{P} \\
            \hline
            Constant & 12.8 & 1.2   & 10.67 & 0.000 \\
            Age      & $-0.12$ & 0.037 & $-3.24$ & 0.003 \\
        \end{tabular}
        \end{center}

        \begin{enumerate}[label=\alph*.]
          \item State H$_0$ and H$_a$ in terms of $\beta$. Define $\beta$ in context.

                \ans{Let $\beta =$ the true slope of the population regression line relating
                maximum running speed to age for adult recreational runners.\\[3pt]
                H$_0$: $\beta = 0$ \quad H$_a$: $\beta \ne 0$}

          \item Identify the test statistic from the output and verify it by calculation.
                Assume LINER conditions are met.

                $t = \dfrac{b}{SE_b} = \dfrac{{\color{keyred}\mathbf{-0.12}}}{{\color{keyred}\mathbf{0.037}}} = \ans{-3.24}$

                \vspace{0.05in}
                $df = n - 2 = \ans{26}$

                \ans{The calculated value $t = -3.24$ matches the value in the output table.}

          \item Interpret the $p$-value in context.

                \ans{Assuming the true slope is 0 (no linear relationship between age and
                maximum running speed), there is a 0.003 probability of obtaining a sample
                slope as extreme as $b = -0.12$ mph per year just by chance.}

          \item State a conclusion at $\alpha = 0.05$.

                \ans{Since $p = 0.003 < 0.05 = \alpha$, we reject H$_0$. There is convincing
                evidence that the true slope of the population regression line relating maximum
                running speed to age for adult recreational runners is not zero. The data support
                the claim that age is linearly related to maximum running speed.}
        \end{enumerate}

  \item \textbf{Absences and Course Grade.} A researcher examines how the number of absences
        ($x$) relates to final course grade ($\%$, $y$) in an introductory statistics course
        (random sample, $n = 35$). Assume LINER conditions are met.

        \begin{center}
        \renewcommand{\arraystretch}{1.4}
        \begin{tabular}{lrrrr}
            \textbf{Predictor} & \textbf{Coef} & \textbf{SE Coef} & \textbf{T} & \textbf{P} \\
            \hline
            Constant  & 91.4 & 2.3  & 39.74 & 0.000 \\
            Absences  & $-2.3$ & 0.54 & $-4.26$ & 0.000 \\
        \end{tabular}
        \end{center}

        \begin{enumerate}[label=\alph*.]
          \item Carry out a complete 4-step significance test at $\alpha = 0.05$.

                \ans{\textbf{Step 1 --- Hypotheses.} Let $\beta =$ the true slope of the
                population regression line relating final course grade to number of absences
                for students in intro statistics. H$_0$: $\beta = 0$; H$_a$: $\beta \ne 0$.\\[3pt]
                \textbf{Step 2 --- Conditions.} LINER conditions are assumed to be met.
                The data come from a random sample.\\[3pt]
                \textbf{Step 3 --- Test Statistic \& $p$-value.}
                $t = -4.26$, $df = 33$, $p \approx 0.000$ (from output).\\[3pt]
                \textbf{Step 4 --- Conclusion.} Since $p \approx 0 < 0.05 = \alpha$, we
                reject H$_0$. There is convincing evidence that the true slope of the
                population regression line relating final course grade to number of absences
                for intro statistics students is not zero.}

          \item Construct a 95\% confidence interval for $\beta$ using $t^* = 2.035$
                ($df = 33$). Show your work.

                \ans{$b \pm t^* \cdot SE_b = -2.3 \pm 2.035(0.54) = -2.3 \pm 1.099
                = (-3.40,\ -1.20)$}

          \item Explain why the confidence interval and the hypothesis test give
                consistent conclusions.

                \ans{Zero is not contained in the 95\% CI $(-3.40,\ -1.20)$. This is
                consistent with rejecting H$_0$: $\beta = 0$ at $\alpha = 0.05$. Both
                procedures agree that $\beta \ne 0$; the interval further indicates that
                the true slope is negative (each additional absence is associated with a
                decrease in course grade).}
        \end{enumerate}

\end{enumerate}

\end{document}
